\documentclass[12pt,UTF8,fontset=none]{ctexbook}
% 导入settings，包含大多数数学书籍常用package
% File settings - applied to all
% 导入第三方库
% \usepackage[fontset=none]{ctex}
\usepackage{amsmath,amssymb,amsthm}
\usepackage{thmtools, thm-restate}
\usepackage[most]{tcolorbox}
\usepackage{enumitem}
\usepackage{geometry}
\usepackage{array}
\usepackage{graphicx}
\usepackage{wrapfig}
\usepackage[table,dvipsnames]{xcolor}
\usepackage{tabularx}
\usepackage{longtable}
\usepackage{float}
\usepackage{dsfont}
\usepackage{mathrsfs}
\usepackage{mathtools}
\usepackage{polynom}
\usepackage{xfrac}
\usepackage{eucal}
\usepackage{titlesec}
\usepackage[version=4]{mhchem}
\usepackage{tikz-cd}
\usepackage{ulem}
\usepackage{verbatim}
\usepackage[makeroom]{cancel}
\usepackage[toc,page]{appendix}
\usepackage{fontspec,xunicode,xltxtra}
\usepackage{caption}
\usepackage[b]{esvect}
\usepackage{thmtools, thm-restate}
\usepackage{pifont}
\usepackage[perpage,symbol*]{footmisc}
\usepackage{xparse} % For defining commands with optional arguments
\usepackage{tikz-3dplot}
\usepackage{pgfplots}
\usepackage{empheq}

% 修改脚注的编号为加圈样式，并且各页单独编号
\DefineFNsymbols{circled}{{\ding{192}}{\ding{193}}{\ding{194}}
{\ding{195}}{\ding{196}}{\ding{197}}{\ding{198}}{\ding{199}}{\ding{200}}{\ding{201}}}
\setfnsymbol{circled}
% \makeatletter
% % 定义 1–20 的圆圈数字
% \newcommand*{\circled@num}[1]{%
%   \ifcase#1\or\ding{172}\or\ding{173}\or\ding{174}\or\ding{175}\or
%         \ding{176}\or\ding{177}\or\ding{178}\or\ding{179}\or\ding{180}\or
%         \ding{181}\or⑪\or⑫\or⑬\or⑭\or⑮\or⑯\or⑰\or⑱\or⑲\or⑳\fi}
% \newcommand*{\circled@nums}{}
% \@for\next:=1,2,...,20\do{%
%   \edef\circled@nums{\circled@nums{\expandafter\circled@num\expandafter{\next}}}
% }
% \DefineFNsymbols*{circled20}{\circled@nums}
% \setfnsymbol{circled20}
% \makeatother

% 自定义颜色
\definecolor{gl}{RGB}{246, 252, 240}
\definecolor{gd}{RGB}{236, 244, 230}
\definecolor{bg}{RGB}{242, 244, 228}
\definecolor{converge_blue}{RGB}{20,76,237}
\definecolor{converge_Blue}{RGB}{55,148,251}
% \definecolor{grid3}{RGB}{200, 200, 200} % 自定义颜色 (RGB 值)
% \tdplotsetmaincoords{60}{105} % 设置观察角度 (θ, φ)

% 选修章节：正文标题、目录、页眉自动加星
\makeatletter
\newcommand{\sectionopt}[2][]{%
  % #1 可选目录短标题，#2 正文标题
  \def\optstar{*}%                                   % 星号
  \def\optmark{\thesection\quad#2\optstar}%           % 页眉格式
  % 1. 目录条目：手动写带*
  \addcontentsline{toc}{section}{%
    \numberline{\thesection}#2\optstar\ifx&#1&\else\space(#1)\fi}%
  % 2. 页眉
  \sectionmark{#2\optstar}%
  % 3. 正文标题
  \refstepcounter{section}%
  \section*{\thesection\quad#2\optstar}%
}
\makeatother

% % 定义字体
\setCJKmainfont[
  BoldFont=STZhongsong,
  ItalicFont=STSong,
  BoldItalicFont=STZhongsong,
  SlantedFont=STSong,
  BoldSlantedFont=STZhongsong]{STSong}  % 普通字体、粗体
\setCJKmonofont[
  BoldFont = simkai.ttf,
  ItalicFont = simkai.ttf,
  BoldItalicFont = simkai.ttf,
  SlantedFont = simkai.ttf,
  BoldSlantedFont = simkai.ttf]{simkai.ttf} % \texttt
\setCJKsansfont[
  BoldFont = simfang.ttf,
  ItalicFont = simfang.ttf,
  BoldItalicFont = simfang.ttf,
  SlantedFont = simfang.ttf,
  BoldSlantedFont = simfang.ttf]{simfang.ttf} % \textsf

% 自制命令
\pgfplotsset{compat=1.18}
\usetikzlibrary{calc,topaths,3d,shadows,positioning,arrows.meta}
\usepgfplotslibrary{colormaps}
\renewcommand{\thesection}{\arabic{chapter}.\arabic{section}}  % 章节使用阿拉伯数字
\renewcommand{\parallel}{\mathrel{/\mskip-4mu/}}  % 平行符号
\renewcommand{\proofname}{\indent\bf 证明}  % 自定义证明标题
\renewcommand{\qedsymbol}{\hfill$\square$}  % 自定义证毕符号
\renewcommand{\Pr}[1]{\mathbb{P}\left(#1\right)}
\newcommand{\e}{\mathrm{e}}  % 自然底数
\newcommand{\dash}{\,–}  % 短折号
\newcommand{\hin}[2]{[\![#1,#2]\!]}  % 连续整数罗列集合 \hin{1}{n}：从1到n的整数，前n个正整数，...
\newcommand{\tong}[1]{\;\overset{#1}{\equiv\joinrel\equiv}\;}  % 同余等号
\newcommand{\di}[1]{\,\mathrm{d}#1}  % 微元d
% \newcommand{\diff}{\,\mathscr{D}}  % 微元d
\newcommand{\diff}{\,\mathrm{d}}  % 微元d
\newcommand{\qu}[2]{\displaystyle\left(#1;#2\right)}  % 开区间
\newcommand{\kon}[2]{\overset{\circ}{U}\left(#1,\,#2\right)}  % 去心邻域
% 局部展开 developpements limites
\newcommand{\oveq}[1]{\overset{#1}{=}}   % equal over
\newcommand{\olim}[1]{\mathit{o}\left(#1\right)}  % petit o
\newcommand{\Olim}[1]{\mathcal{O}\left(#1\right)}  % grand O
\newcommand{\Tlim}[1]{\tau\left(#1\right)}  % grand theta
\newcommand{\eqlim}[1]{\overset{#1}{\sim}}  % equivalence
\newcommand{\vect}[1]{\left\langle #1 \right\rangle}  % 生成空间 generated space

\newcommand{\arccot}{\operatorname{arccot}}  % 反余弦函数
\newcommand{\dlim}[1]{^{\color{gray}\prime}#1}  % 数字分隔符
\newcommand{\lian}[1]{  % 极限符号
    \underset{#1}{\operatorname{lian}\,}
}
\newcommand{\jie}[3]{\displaystyle\left( \,#1 \,\middle\vert \,#2\,\middle\vert \,#3 \,\right)}  % 内积
\newcommand{\nji}[2]{\displaystyle\left( \,#1 \,\middle\vert \,#2 \,\right)}  % 内积
\newcommand{\supp}[1]{\underset{#1}{\text{\texttt{上}}}}  % 上确界
\newcommand{\inff}[1]{\underset{#1}{\text{\texttt{下}}}}  % 下确界
\newcommand{\maxx}[1]{\underset{#1}{\text{\texttt{大}}}}  % 最大值
\newcommand{\minn}[1]{\underset{#1}{\text{\texttt{小}}}}  % 最小值
\newcommand{\dangle}{  % 平面角符号
    \mathord{
        \tikz[baseline=(current bounding box.south), 
          trim left=0pt, trim right=1.6em] \draw (0.8em,0ex) -- (0.3em, 0ex) -- (.6em, 1.5ex) -- (.8em, 1.5ex) -- (.5em, 0ex) -- cycle;
    }
}
\newcommand{\xangle}{  % 平面角符号
    \mathord{
        \tikz[baseline=(current bounding box.south), 
          trim left=0pt, trim right=1.6em] \draw (0.8em,1.5ex) -- (0.3em, 0ex) -- (.64em, 0ex) -- (.8em, .36ex) -- (.42em, .36ex) -- cycle;
    }
}
\newcommand{\bu}{  % 补集符号
  \mathord{%
    \tikz[baseline=-0.1ex, trim left=0pt, trim right=0.8ex] {%
      \fill (0,0) rectangle (0.8ex,0.8ex);
      \fill[white] (0.4ex,0.4ex) circle (0.3ex);
    }%
  }%
}
\newcommand{\inv}[1]{#1^{(-1)}}  % 映射的逆映射
\newcommand{\rectbx}{  % 长方形符号
    \mathord{
        \tikz[baseline=(current bounding box.south), 
          trim left=0pt, trim right=0.7em] \draw (0,.1ex) -- (.5em,.1ex) -- (.5em,0.8em) -- (0em,0.8em) -- cycle;
    }
}
\newcommand{\tr}{  % 矩阵转置符号 A^{\tr} 
    \mathord{
        \begin{tikzpicture}[baseline=-0.2em, line width=0.5pt]
        \draw (-0.15em, 0.15em) -- (0.06em, -0.06em);
        \draw (45:0.15em) arc[start angle=45, end angle=225, radius=0.15em];
    \end{tikzpicture}
    }
}
% \newcommand{\ball}{
%     \mathord{%
%         \tikz[baseline=0.2em, trim left=0pt, trim right=0.9em] {%
%             \draw (0.45em,0.45em) circle (0.45em);
%             \draw (0em,0.45em) arc [start angle=180, end angle=360, x radius=0.45em, y radius=0.15em];
%             \draw (0.45em,0.9em) arc [start angle=90, end angle=270, x radius=0.15em, y radius=0.45em];
%             \draw[dash pattern={on 1.4pt off 1.2pt}] (0em,0.45em) arc [start angle=180, end angle=0, x radius=0.45em, y radius=0.15em];
%             \draw[dash pattern={on 1.4pt off 1.2pt}] (0.45em,0.9em) arc [start angle=90, end angle=-90, x radius=0.15em, y radius=0.45em];
%         }%
%     }%
% }
\newcommand{\arcangle}{\mathord{\mathpalette\doarcangle\relax}}  % 带弧的角度符号 - 交角
\newcommand{\doarcangle}[2]{  % 
    \hbox{%
        \sbox0{$#1B$}%
        \sbox2{$#1<$}%
        \raisebox{\dimexpr\dp0+(\ht0-\ht2)/2}{%
            $#1<\mspace{-9mu}\mathrel{)}\mspace{2mu}$%
        }%
    }%
}
\newcommand{\parasbx}{  % 平行四边形符号
    \mathord{
        \tikz[baseline=(current bounding box.south), 
          trim left=0pt, trim right=1em] \draw (0,.1ex) -- (.8em,.1ex) -- (1em,1.6ex) -- (.2em,1.6ex) -- cycle;
    }
}
\newcommand{\widearc}[1]{  % 可伸缩圆弧符号
    \tikz[baseline=(wideArcAnchor.base)]{
        \node[inner sep=0] (wideArcAnchor) {$#1$}; 
        \coordinate (wideArcAnchorA) at ($(wideArcAnchor.north west) + (0.15em,0.1em)$);
        \coordinate (wideArcAnchorB) at ($(wideArcAnchor.north east) + (0.0em,0.1em)$);
        \draw[line width=0.1ex,line cap=round,out=45,in=135] (wideArcAnchorA) to (wideArcAnchorB);
    }
}
\newcommand{\bino}[2]{\mathrm{C}_{#1}^{#2}}
\NewDocumentCommand{\pilo}{m o}{%
    \mathrm{P}_{#1}\IfValueT{#2}{^{#2}}
}
\NewDocumentCommand{\bmat}{m o}{%
    \left[\,#1\,\right]\IfValueT{#2}{_{#2}}
}
\newcommand{\fanshe}[1]{\text{反}_{#1}}
\newcommand{\touying}[1]{\text{投}_{#1}}
\newcommand{\xuan}[1]{\text{旋}_{#1}}
% 定义选学章节命令：目录和标题同步加星号
\newcommand{\chapterext}[1]{%
  \chapter[#1\textsuperscript{*}]{#1\textsuperscript{*}}%
}
\newcommand{\sectionext}[1]{%
  \section[#1\textsuperscript{*}]{#1\textsuperscript{*}}%
}
\newcommand{\subsectionext}[1]{%
  \subsection[#1\textsuperscript{*}]{#1\textsuperscript{*}}%
}
\newcommand{\jii}{\int\!\!\!\!\!\int}
\newcommand{\jiil}{\mathop{\int\!\!\!\!\!\int}\limits}
\newcommand{\jiii}{\int\!\!\!\!\!\int\!\!\!\!\!\int}
\newcommand{\jiiil}{\mathop{\int\!\!\!\!\!\int\!\!\!\!\!\int}\limits}

% 定义、定理、证明等块环境
\theoremstyle{definition}
\newtheorem{df}{定义}[section] 
\newtheorem*{po}{公理}
\newtheorem{pp}{命题}[section]
\newtheorem{tm}{定理}[section]
\newtheorem{cor}{推论}[pp]
\newtheorem{ex}{例子}[section]
\newtheorem{et}{例题}[section]
\newtheorem*{ex*}{例子}
\newtheorem*{so}{解答}
\theoremstyle{plain}
\newtheorem{sk}{思考}[section]
\newtheorem{xt}{习题}[section]
\renewenvironment{proof}{\paragraph{\textbf{证明：}}}{\hfill$\square$}

\setlength{\intextsep}{2pt}%
\setlength{\columnsep}{2pt}%
% 列举环境
\setlist{label=\textbullet}
% 列举环境行间距
\setenumerate[1]{itemsep=0pt,partopsep=0pt,parsep=0pt,topsep=0pt}
\setenumerate[2]{itemsep=0pt,partopsep=0pt,parsep=\parskip,topsep=0pt}
\setitemize[1]{itemsep=0pt,partopsep=0pt,parsep=0pt,topsep=0pt}
\setdescription{itemsep=0pt,partopsep=0pt,parsep=0pt,topsep=0pt}
% 章节间距
\setlength\parskip{8pt}
% 文本框间距
\setlength{\fboxsep}{12pt}
% 章节字体大小
\titleformat{\section}{\zihao{-2}\bfseries}{ \thesection }{16pt}{}

% 封面
\title{\zihao{0} \bfseries XXXX}
\author{\zihao{2} \texttt{XXX}}
% \date{\bfseries\today}
\date{}
% 正文
\begin{document}
\maketitle
\tableofcontents
\newpage

\chapter{欧拉定理}

\section{从一个数数问题开始}

我们先来看几个常见的多面体。多面体是由平面多边形围成的立体，例如正方体、四棱锥和三棱柱。对每一个多面体，数一数它的顶点数、棱数和面数，然后计算顶点数减去棱数再加上面数，也就是 $V-E+F$。

正方体有 $8$ 个顶点、$12$ 条棱和 $6$ 个面，于是 $V-E+F=8-12+6=2$。四棱锥有 $5$ 个顶点、$8$ 条棱和 $5$ 个面，$5-8+5=2$。三棱柱有 $6$ 个顶点、$9$ 条棱和 $5$ 个面，$6-9+5=2$。

% 插图：一张组合图，展示正方体、四棱锥、三棱柱的立体图，每个多面体旁边用不同颜色标出顶点、棱、面，并注明数量。
% 插图要求：画面清晰，三个多面体大小适中，顶点可用小圆点表示，棱用实线，面用浅色填充。

这些结果都是 $2$。如果仅仅三个例子还不能说明问题，不妨再看几个更复杂的凸多面体。正十二面体有 $20$ 个顶点、$30$ 条棱和 $12$ 个面，$20-30+12=2$。正二十面体有 $12$ 个顶点、$30$ 条棱和 $20$ 个面，$12-30+20=2$。

% 插图：一张表格，表格中列出正方体、四棱锥、三棱柱、正十二面体、正二十面体，以及各自的 V、E、F 和 V-E+F 的值。
% 插图要求：表格用横线分隔，数值居中，最后一行可加粗显示结果均为 2。

所有试过的凸多面体，无论形状如何，这个值总是 $2$。这显然不是巧合，而是一个普遍的规律。那么，为什么偏偏是 $2$？为什么不是 $0$ 或 $3$？如果换一个有洞的多面体，比如像甜甜圈那样的形状，结果还会是 $2$ 吗？

这个奇怪的现象早在十八世纪就被欧拉注意到了。他发现，对于任何凸多面体，顶点数 $V$、棱数 $E$、面数 $F$ 之间恒有
\[
V-E+F=2.
\]
这就是后来所称的欧拉定理。然而，仅仅知道这个公式还不够。它的成立依赖于多面体的哪些性质？为什么数值恰好是 $2$？这些问题的答案隐藏在比多面体本身更基本的图形性质之中。接下来，我们将从一个更简单的对象——平面上的图——出发，逐步揭开这个秘密。

\section{先解决一个更简单的问题：平面图的欧拉公式}

多面体是立体的，直接研究它并不方便。一个自然的想法是，把多面体表面摊平到平面上，把立体问题转化为平面问题。在这样做之前，先来看一类更基本的对象——平面图。

所谓平面图，就是画在平面上的一张图：由一些顶点和连接顶点的边组成，要求任意两条边除端点外互不相交。平面图中的边不一定是直线段，只要彼此不交叉就可以。图被它的边分割成若干连通的区域，每个这样的区域称为一个面。需要特别注意的是，平面图还有一个外部面，它位于图的边缘之外，一直延伸到无穷远处。在数面数的时候，这个外部面也要算进去。

% 插图：一张简单平面图，包含内部若干区域和一个外部面。
% 插图要求：图中顶点用黑点标出，边用曲线表示，内部面用不同浅色填充，外部面用最浅色填充并标注“外部面”。

我们先从最简单的平面图——树——开始。树是一种没有圈的连通图。所谓圈，是指一条首尾相接、中途不经过重复顶点的闭合路径。树没有圈，因此它不会把平面分割成多个部分。一棵树作为平面图时，整个平面仍然是一个连通的区域，所以它只有一个面，即外部面。于是对于树，$F=1$。

接下来证明树还满足 $E=V-1$。这个事实可以用归纳法来证明。当树只有一个顶点时，它没有边，结论显然成立。假设结论对所有顶点数较少的树成立，现在来看一棵有 $V$ 个顶点的树。树一定有叶子，也就是只连接一条边的顶点。如果从任意一个顶点出发沿着边走，因为图是有限的，又因为没有圈而不会走回已经经过的顶点，所以最终一定会停在一个度数为 $1$ 的顶点上，这就是叶子。删去这个叶子以及它所连接的唯一一条边，剩下的图形仍然是一棵树，顶点数变为 $V-1$，边数变为 $E-1$。由归纳假设，$E-1=(V-1)-1$，因此 $E=V-1$。于是对于树，
\[
V-E+F=V-(V-1)+1=2.
\]
这说明任何树都满足欧拉公式。

% 插图：一棵树的结构图，其中一个叶子被圆圈标出，旁边用虚线画出删去该叶子及边后的效果。
% 插图要求：左边是原树，右边是删去叶子后的树，虚线示意删除动作。

现在考虑一般的连通平面图。如果图中含有圈，我们可以在圈上选取一条边。圈是一条简单闭曲线，它把平面分成内外两部分，因此这条被选取的边两侧是两个不同的面。删去这条边后，圈上出现一个缺口，内外两个区域通过缺口连通，于是这两个面合并成一个面。其他面不受影响。所以删去圈上一条边后，顶点数不变，边数减少 $1$，面数也减少 $1$，而图仍然连通。此时 $V-E+F$ 的值保持不变。

% 插图：一个带圈的平面图，图中一条边用粗线标出，旁边用虚线展示删去该边后面数减少的效果。
% 插图要求：左图用两种颜色区分被标边两侧的两个不同面，右图显示删去该边后两色区域合并。

这样一来，只要图中有圈，就可以反复删去圈上的边。每删去一条这样的边，$V-E+F$ 的值都不变，而边数不断减少。这个过程最终会得到一个连通且没有圈的图，也就是树。我们已经知道树满足 $V-E+F=2$，所以原图也满足 $V-E+F=2$。

这个结论可以更正式地表述为：任意连通平面图都有
\[
V-E+F=2.
\]
它被称为平面图的欧拉公式。值得注意的是，这个证明完全在平面内完成，没有涉及任何立体几何或拓扑变形。它只依赖于一个简单而核心的事实：树不产生圈，而每添加一条边，若形成圈，就会增加一个面，同时减少一条边，使得 $V-E+F$ 保持不变。这正是公式中那个常数 $2$ 的来源。

% 插图：一组三格漫画风格的图，展示从一般连通平面图通过不断删去圈上边最终变成树的过程。
% 插图要求：三张小图从左到右排列，第一张是有圈的平面图，第二张删去两条边后面数减少，第三张变成树，所有图中标注 V-E+F 的值始终为 2。

\section{从多面体到平面图}

平面图的欧拉公式已经证明完毕，但我们的目标是多面体。一个自然的想法是，把多面体的表面摊到平面上，变成一张平面图，这样就能直接使用前面的结论。这个“摊平”的过程当然不能真的用剪刀和胶水，而需要在想象中完成。可以这样设想：多面体的表面是一层极薄的橡皮膜，可以任意拉伸、压缩、弯曲，只要不撕破，也不把不同的点粘在一起。先把它吹成一个球面，再在球面上剪开一个小洞，把剩下的部分压平到桌面上。经过这样一番操作，多面体上的顶点、棱、面只是挪动了位置，它们的连接关系没有发生任何改变。

这种能拉伸、能压缩，但不撕裂、不粘合的变化，在数学中有一个专门的名字，叫做同胚。更准确地说，如果两个图形能够通过连续的伸缩、弯曲互相变成对方，并且整个过程不撕破也不粘连，那么这两个图形就是同胚的。比如，一个气球表面和一个球面是同胚的，因为气球可以吹成球面，也可以缩成各种形状的球面；一个橡皮圆盘和一张平面也是同胚的，因为圆盘边缘可以向外拉伸，最终铺满整个平面。同胚所关心的不是大小、弯曲程度，甚至不是直线或曲线，而只是点与点之间的连接关系。因此，如果多面体表面同胚于球面，而球面去掉一个点后又同胚于平面，那么把多面体表面看成平面图就是合理的，顶点数、棱数、面数都不会改变。

剩下的问题是，凸多面体的表面是否真的同胚于球面。这件事在直觉上似乎很明显，但要给出一个严格的说明，还需要一点细致的推理。

首先需要一个位于多面体内部的点。设多面体是凸的，也就是说，多面体上任意两点之间的线段全都落在多面体内。在这个前提下，把多面体所有顶点的坐标取算术平均，得到的点记作 $O$。由于 $O$ 是各顶点的凸组合，而凸多面体对凸组合是封闭的，所以 $O$ 一定在多面体内。它还不能在表面上。如果 $O$ 在表面上，那么它必定落在某个面所在的平面上。由于多面体是凸的，所有顶点都在这个平面的同一侧或就在平面上。但 $O$ 是所有顶点的平均，如果某个顶点不在该平面上，平均之后就会偏向那一侧，不可能恰好落在平面上。因此所有顶点都必须在那个平面上，这样多面体就退化成了一个平面图形，这与我们讨论的立体多面体矛盾。所以 $O$ 是真正的内部点，也就是说，存在一个以 $O$ 为球心的小球完全包含在多面体内部。

% 插图：一个凸多面体的剖面图，内部标出顶点平均点 O，并用虚线画出一个以 O 为球心的小球完全包含在多面体内部。
% 插图要求：剖面图显示多面体内部，O 点明显，小球与表面保持一定距离，旁边标注“内部点 O”。

现在从 $O$ 出发作任意一条射线。由于 $O$ 是内部点，射线一开始的一小段仍在多面体内；又因为多面体有界，沿射线走得足够远就会离开多面体。从内部到外部连续变化，必然第一次穿过表面，所以射线与表面至少有一个交点。如果有两个交点 $A$ 和 $B$，且从 $O$ 出发先经过 $A$ 再经过 $B$，那么由凸性，线段 $AB$ 整个在多面体内。又因为 $O$ 是内部点，存在一个以 $O$ 为球心的小球完全在多面体内。把小球中的每个点与 $B$ 相连，这些线段都在多面体内，它们在 $A$ 附近会围成一个小区域，把 $A$ 完全包在多面体内部。这与 $A$ 是表面点矛盾。因此从 $O$ 出发的射线与多面体表面恰好有一个交点。

于是，以 $O$ 为球心作一个足够大的球面，把多面体完全包含在内。对多面体表面上的每一点 $X$，从 $O$ 出发经过 $X$ 的射线与球面交于唯一一点 $X'$。这个对应将多面体表面的每一点都映到球面上，并且不同的点映到不同的点，反过来球面上每一点也对应多面体表面上的一个点。因此这是多面体表面到球面的一一对应，而且对应是连续的。这就说明凸多面体表面与球面同胚。

% 插图：一个凸多面体内部取点 O，外部画一个大球面，从 O 出发的几条射线将多面体表面上的点投影到球面上。
% 插图要求：多面体用实线表示，球面用虚线表示，几条射线从 O 发出，穿过多面体表面并延伸到球面，对应点用相同字母标记。

接下来的任务是把球面摊成平面。在多面体表面上任意选一个面，记作 $F_0$。它在球面上对应一个区域。在这个区域内取一点 $P$，例如取 $F_0$ 的中心在球面上的对应点。现在把球面从点 $P$ 进行球极投影。球极投影的做法是：把球面看成一个透明球，在点 $P$ 处放一个点光源，让球面上的图形投影到点 $P$ 的对径点处的切平面上。除了点 $P$ 本身之外，球面上每一点都唯一对应切平面上的一点，反过来也成立。因此球面去掉点 $P$ 后与整个平面是一一对应的。更一般地，球面去掉一个点后与平面同胚。

% 插图：球极投影示意图。球面上标出点 P 和对径点 Q，平面为过 Q 的切平面，从 P 发出的几条光线将球面上的点投影到平面上。
% 插图要求：球面用实线，切平面用虚线，点 P 在球面顶部，点 Q 在底部，几条从 P 出发的直线穿过球面到达平面，投影点清晰。

经过球极投影，多面体表面的图形被完整地铺到了平面上。原来被去掉的面 $F_0$ 所对应的区域包含点 $P$，投影后变成了整个平面的外部无界区域，其余的面则成为平面上的有界区域。顶点和棱也一一对应地落在平面上。这样得到的图形满足平面图的要求：顶点之间由边连接，任意两条边除端点外不相交，整个平面被分成若干个面。顶点数、棱数、面数都与原来的多面体相同。

至此，凸多面体表面通过同胚变形成为平面图的过程已经清楚。既然同胚保持顶点、边、面的连接关系，而平面图的欧拉公式已经证明，那么对于任意凸多面体，立即得到
\[
V-E+F=2.
\]
这就是欧拉定理。值得注意的是，整个过程中真正起作用的不是多面体的具体形状，而是它表面与球面的同胚关系。正是这种同胚关系保证了平面图的方法能够适用。

% 插图：正方体到平面图的同胚变形示意图，分三格展示：正方体表面投影到球面，球面去掉一个面后球极投影到平面，平面图标注出外部面和内部面。
% 插图要求：三格从左到右排列，第一格是正方体与球面投影关系，第二格是球极投影，第三格是摊平后的平面图。平面图中外部面用浅色阴影表示，内部面用不同颜色区分。

\section{从欧拉示性数到拓扑学}

回顾整个证明过程，欧拉定理的成立依赖于两个层次的事实。第一层是图论中的平面图欧拉公式：任意连通平面图都满足 $V-E+F=2$。这个结论完全可以在平面内用归纳法证明，它的核心在于树没有圈，而圈的增加会同时增加面和边，从而保持 $V-E+F$ 不变。第二层是拓扑中的同胚概念：凸多面体表面与球面同胚，球面去掉一个点后与平面同胚，因此多面体表面可以转化为平面图，顶点数、棱数、面数均保持不变。这两个层次结合起来，便得到了欧拉定理。

进一步看，$V-E+F$ 这个量对于同胚的曲面来说保持不变。对于球面，它是 $2$；对于环面，它是 $0$；对于双环面，它是 $-2$。这表明 $V-E+F$ 并不依赖于具体的多面体剖分，而只依赖于曲面本身的形状。这样的量称为拓扑不变量。$V-E+F$ 被特别地称为欧拉示性数。

拓扑学所关心的正是同胚变换下保持不变的性质。欧拉示性数就是其中最简单、最直观的一个例子。通过计算一个曲面的欧拉示性数，可以判断它是否与另一个曲面同胚。球面与环面不同胚，正是因为它们的欧拉示性数不同。曲面上的洞的个数与欧拉示性数之间存在确定的关系：每增加一个洞，欧拉示性数减少 $2$。这个关系在初等几何中只能通过具体例子观察，但要在一般情形下严格证明，需要更系统的工具。

拓扑学发展出了一套完整的理论来处理这类问题。其中，单纯同调是一种基本的方法。它将曲面剖分成三角形或其他单纯形，然后通过代数手段计算欧拉示性数，并证明它确实是拓扑不变量。这套方法推广了平面图中数顶点、边、面的思想，使它适用于更高维的空间。

欧拉定理最初只是一个关于多面体的计数公式，但它所揭示的规律远远超出了多面体本身。从平面图的归纳证明，到同胚变形的直观理解，再到欧拉示性数和拓扑不变量的建立，这个简单的公式将初等几何与一个广阔的数学领域连接起来。通过对这个公式的考察，我们已经初步看到了拓扑学的基本问题和基本方法。

\end{document}